% GENERATED FILE. Edit canonical lessons, then run npm run generate:books.
% canonical-source-hash: fnv1a64:a0ce032695859006
% canonical-lessons: SA-C373-anumatih, SA-C373-nisedhati, SA-C373-pratibandhah, SA-C373-apavadah, SA-C373-anivaryah

\chapter{Permission First}
\label{ch:sa-permission-first}
\hlchaptermodality{373}

\begin{chapteropening}
\textbf{By the end of this chapter:} \emph{I can say what is allowed, what is forbidden and what is required.}

Complete the last lesson of chapter 373: I can say what is allowed, what is forbidden and what is required.
\end{chapteropening}

\section[\texorpdfstring{\sk{अनुमतिः} (anumatiḥ)}{अनुमतिः (anumatiḥ)}]{\sk{अनुमतिः} (anumatiḥ) — permission}

\label{lesson:SA-C373-anumatih}



\begin{quote}
Before the new one: say the Sanskrit for holds up, then the Sanskrit for removes.
\end{quote}

\subsection*{You'll want to know: \sk{अनुमतिः}}
\textbf{\sk{अनुमतिः}} — \emph{anumatiḥ} — \textquotedblleft{}permission\textquotedblright{}.

\begin{quote}
\sk{यदि} \sk{पिता} \sk{अनुमतिं} \sk{ददाति}, \sk{तर्हि} \sk{अहं} \sk{गमिष्यामि।}
\end{quote}

\begin{quote}
\emph{yadi pitā anumatiṃ dadāti, tarhi ahaṃ gamiṣyāmi.}
\end{quote}

\begin{quote}
If father gives permission, I will go.
\end{quote}

A feminine noun.

\begin{cousinweb}
\textbf{\sk{अनुमतिः}} (\emph{anumatiḥ}) is built on \emph{anu}, \textquotedblleft{}after\textquotedblright{}, and the root √man, \textquotedblleft{}to think\textquotedblright{}; MW gives \emph{anu-mati} \textquotedblleft{}assent, permission, approbation\textquotedblright{} (MW).
\end{cousinweb}

\begin{grammarlens}[title={What you've built}]
One more word for saying what happens if.
\end{grammarlens}

\subsection*{Your turn}
\begin{itemize}
\raggedright
  \item \emph{Say it:} \emph{anumatiḥ}
  \item \emph{Say it:} \emph{anumatiḥ}, once more
  \item \emph{Say it:} \emph{yadi pitā anumatiṃ dadāti, tarhi ahaṃ gamiṣyāmi.} — If father gives permission, I will go.
  \item \emph{Recall:} say the Sanskrit for holds up, then the Sanskrit for removes, then say \emph{anumatiḥ} again
\end{itemize}

\begin{tcolorbox}[breakable,colback=teal!4,colframe=teal!35!black,title={Before you move on}]
What does \sk{अनुमतिः} mean? (\textquotedblleft{}Permission\textquotedblright{}.) Say it once more.
\end{tcolorbox}
\section[\texorpdfstring{\sk{निषेधति} (niṣedhati)}{निषेधति (niṣedhati)}]{\sk{निषेधति} (niṣedhati) — forbids}

\label{lesson:SA-C373-nisedhati}



\begin{quote}
Before the new one: say the Sanskrit for removes, then the Sanskrit for permission.
\end{quote}

\subsection*{You'll want to know: \sk{निषेधति}}
\textbf{\sk{निषेधति}} — \emph{niṣedhati} — \textquotedblleft{}forbids\textquotedblright{}.

\begin{quote}
\sk{यदि} \sk{वैद्यः} \sk{निषेधति}, \sk{तर्हि} \sk{सः} \sk{मिष्टान्नं} \sk{न} \sk{खादति।}
\end{quote}

\begin{quote}
\emph{yadi vaidyaḥ niṣedhati, tarhi saḥ miṣṭānnaṃ na khādati.}
\end{quote}

\begin{quote}
If the doctor forbids it, he does not eat sweets.
\end{quote}

Past (he …): \sk{सः} \sk{मिष्टान्नं} \sk{निषिद्धवान्}, \emph{saḥ miṣṭānnaṃ niṣiddhavān}.

\begin{cousinweb}
\textbf{\sk{निषेधति}} (\emph{niṣedhati}) is built on \emph{ni} and the root √sidh, \textquotedblleft{}to keep off\textquotedblright{}; MW gives \emph{ni-√ṣidh} \textquotedblleft{}to keep off, to prevent, to prohibit\textquotedblright{} and \emph{ni-ṣiddha} \textquotedblleft{}forbidden\textquotedblright{} (MW).
\end{cousinweb}

\begin{grammarlens}[title={What you've built}]
One more word for saying what happens if.
\end{grammarlens}

\subsection*{Your turn}
\begin{itemize}
\raggedright
  \item \emph{Say it:} \emph{niṣedhati}
  \item \emph{Say it:} \emph{niṣedhati}, once more
  \item \emph{Say it:} \emph{yadi vaidyaḥ niṣedhati, tarhi saḥ miṣṭānnaṃ na khādati.} — If the doctor forbids it, he does not eat sweets.
  \item \emph{Recall:} say the Sanskrit for removes, then the Sanskrit for permission, then say \emph{niṣedhati} again
\end{itemize}

\begin{tcolorbox}[breakable,colback=teal!4,colframe=teal!35!black,title={Before you move on}]
What does \sk{निषेधति} mean? (\textquotedblleft{}Forbids\textquotedblright{}.) Say it once more.
\end{tcolorbox}
\section[\texorpdfstring{\sk{प्रतिबन्धः} (pratibandhaḥ)}{प्रतिबन्धः (pratibandhaḥ)}]{\sk{प्रतिबन्धः} (pratibandhaḥ) — a restriction; an impediment}

\label{lesson:SA-C373-pratibandhah}



\begin{quote}
Before the new one: say the Sanskrit for permission, then the Sanskrit for forbids.
\end{quote}

\subsection*{You'll want to know: \sk{प्रतिबन्धः}}
\textbf{\sk{प्रतिबन्धः}} — \emph{pratibandhaḥ} — \textquotedblleft{}a restriction; an impediment\textquotedblright{}.

\begin{quote}
\sk{यदि} \sk{प्रतिबन्धः} \sk{अस्ति}, \sk{तर्हि} \sk{कः} \sk{अपि} \sk{न} \sk{प्रविशतु।}
\end{quote}

\begin{quote}
\emph{yadi pratibandhaḥ asti, tarhi kaḥ api na praviśatu.}
\end{quote}

\begin{quote}
If there is a restriction, let no one go in.
\end{quote}

A masculine noun.

\begin{cousinweb}
\textbf{\sk{प्रतिबन्धः}} (\emph{pratibandhaḥ}) is built on \emph{prati} and the root √bandh, \textquotedblleft{}to bind\textquotedblright{}; MW gives \emph{prati-bandha} \textquotedblleft{}an obstacle, a hindrance, an impediment\textquotedblright{}, and \textquotedblleft{}a restriction\textquotedblright{} is how modern Sanskrit uses it (MW).
\end{cousinweb}

\begin{grammarlens}[title={What you've built}]
One more word for saying what happens if.
\end{grammarlens}

\subsection*{Your turn}
\begin{itemize}
\raggedright
  \item \emph{Say it:} \emph{pratibandhaḥ}
  \item \emph{Say it:} \emph{pratibandhaḥ}, once more
  \item \emph{Say it:} \emph{yadi pratibandhaḥ asti, tarhi kaḥ api na praviśatu.} — If there is a restriction, let no one go in.
  \item \emph{Recall:} say the Sanskrit for permission, then the Sanskrit for forbids, then say \emph{pratibandhaḥ} again
\end{itemize}

\begin{tcolorbox}[breakable,colback=teal!4,colframe=teal!35!black,title={Before you move on}]
What does \sk{प्रतिबन्धः} mean? (\textquotedblleft{}A restriction; an impediment\textquotedblright{}.) Say it once more.
\end{tcolorbox}
\section[\texorpdfstring{\sk{अपवादः} (apavādaḥ)}{अपवादः (apavādaḥ)}]{\sk{अपवादः} (apavādaḥ) — an exception}

\label{lesson:SA-C373-apavadah}



\begin{quote}
Before the new one: say the Sanskrit for forbids, then the Sanskrit for a restriction.
\end{quote}

\subsection*{You'll want to know: \sk{अपवादः}}
\textbf{\sk{अपवादः}} — \emph{apavādaḥ} — \textquotedblleft{}an exception\textquotedblright{}.

\begin{quote}
\sk{यदि} \sk{नियमः} \sk{अस्ति}, \sk{तर्हि} \sk{अपवादः} \sk{अपि} \sk{भवति।}
\end{quote}

\begin{quote}
\emph{yadi niyamaḥ asti, tarhi apavādaḥ api bhavati.}
\end{quote}

\begin{quote}
If there is a rule, there is also an exception.
\end{quote}

A masculine noun.

\begin{cousinweb}
\textbf{\sk{अपवादः}} (\emph{apavādaḥ}) is built on \emph{apa}, \textquotedblleft{}away\textquotedblright{}, and the root √vad, \textquotedblleft{}to speak\textquotedblright{}; MW gives \emph{apa-vāda} \textquotedblleft{}a special rule setting aside a general one, an exception\textquotedblright{} (MW).
\end{cousinweb}

\begin{grammarlens}[title={What you've built}]
One more word for saying what happens if.
\end{grammarlens}

\subsection*{Your turn}
\begin{itemize}
\raggedright
  \item \emph{Say it:} \emph{apavādaḥ}
  \item \emph{Say it:} \emph{apavādaḥ}, once more
  \item \emph{Say it:} \emph{yadi niyamaḥ asti, tarhi apavādaḥ api bhavati.} — If there is a rule, there is also an exception.
  \item \emph{Recall:} say the Sanskrit for forbids, then the Sanskrit for a restriction, then say \emph{apavādaḥ} again
\end{itemize}

\begin{tcolorbox}[breakable,colback=teal!4,colframe=teal!35!black,title={Before you move on}]
What does \sk{अपवादः} mean? (\textquotedblleft{}An exception\textquotedblright{}.) Say it once more.
\end{tcolorbox}
\section[\texorpdfstring{\sk{अनिवार्यः} (anivāryaḥ)}{अनिवार्यः (anivāryaḥ)}]{\sk{अनिवार्यः} (anivāryaḥ) — compulsory; unavoidable}

\label{lesson:SA-C373-anivaryah}



\begin{quote}
Before the new one: say the Sanskrit for a restriction, then the Sanskrit for an exception.
\end{quote}

\subsection*{You'll want to know: \sk{अनिवार्यः}}
\textbf{\sk{अनिवार्यः}} — \emph{anivāryaḥ} — \textquotedblleft{}compulsory; unavoidable\textquotedblright{}.

\begin{quote}
\sk{यदि} \sk{परीक्षा} \sk{अनिवार्या} \sk{अस्ति}, \sk{तर्हि} \sk{अहं} \sk{पठिष्यामि।}
\end{quote}

\begin{quote}
\emph{yadi parīkṣā anivāryā asti, tarhi ahaṃ paṭhiṣyāmi.}
\end{quote}

\begin{quote}
If the exam is compulsory, I will study.
\end{quote}

\begin{cousinweb}
\textbf{\sk{अनिवार्यः}} (\emph{anivāryaḥ}) is built from \emph{a-}, \textquotedblleft{}not\textquotedblright{}, and \emph{nivārya}, \textquotedblleft{}to be warded off\textquotedblright{}; MW gives \emph{a-nivārya} \textquotedblleft{}not to be warded off, unavoidable\textquotedblright{}, and \textquotedblleft{}compulsory\textquotedblright{} is its modern sense (MW).
\end{cousinweb}

\begin{grammarlens}[title={What you've built}]
One more word for saying what happens if.
\end{grammarlens}

\subsection*{Your turn}
\begin{itemize}
\raggedright
  \item \emph{Say it:} \emph{anivāryaḥ}
  \item \emph{Say it:} \emph{anivāryaḥ}, once more
  \item \emph{Say it:} \emph{yadi parīkṣā anivāryā asti, tarhi ahaṃ paṭhiṣyāmi.} — If the exam is compulsory, I will study.
  \item \emph{Recall:} say the Sanskrit for a restriction, then the Sanskrit for an exception, then say \emph{anivāryaḥ} again
\end{itemize}

\begin{tcolorbox}[breakable,colback=teal!4,colframe=teal!35!black,title={Before you move on}]
What does \sk{अनिवार्यः} mean? (\textquotedblleft{}Compulsory; unavoidable\textquotedblright{}.) Say it once more.
\end{tcolorbox}
